\section{The Labor Market Foundation: Twenty Papers by David Autor}
\label{app:autor}
%%%%%%%%%%%%%%%%%%%%%%%%%%%%%%%%%%%%%%%%%%%%%%%%%%%%%%%%%%%%%%%%%%%%%%%%%%%%%%%

This appendix integrates twenty of the most influential papers by David Autor
(MIT), whom \emph{The Economist} called ``The academic voice of the American worker.''
His work spans technological change, globalization, job polarization, and
labor market adjustment. Together with Heckman (human capital formation),
Autor provides the \emph{labor economics} pillar of the $(C, \Psi)$ framework,
focusing on how technology and trade reshape the task structure of work.

\subsection{Overview: Autor's Research in Framework Terms}

Autor's central contribution is the \emph{task-based approach} to labor markets:
\begin{equation}
  Y = F(L_{routine}, L_{abstract}, L_{manual}, K, \Psi_{tech})
\end{equation}
where labor is differentiated by task type, and technology ($\Psi_{tech}$)
affects different tasks differently:
\begin{itemize}
  \item \textbf{Routine tasks:} Substitutable by computers ($C^{tech,routine} < 0$)
  \item \textbf{Abstract tasks:} Complemented by computers ($C^{tech,abstract} > 0$)
  \item \textbf{Manual tasks:} Largely unaffected (flexible, non-codifiable)
\end{itemize}

\begin{center}
\renewcommand{\arraystretch}{1.2}
\small
\begin{tabular}{ll}
\hline
\textbf{Framework Element} & \textbf{Autor's Contribution} \\
\hline
$\Psi_{tech}$ & Routine-biased technical change \\
Task taxonomy & Routine, abstract, manual \\
$C^{tech,task}$ varies & Technology affects tasks differently \\
Job polarization & Middle-skill jobs disappear \\
China shock & Trade as $\Psi$ shift, local effects \\
Slow adjustment & Labor markets adjust incompletely \\
\hline
\end{tabular}
\end{center}

\subsection{Part I: The Task Approach to Technological Change (6 Papers)}

These papers established the task-based framework for understanding how
technology transforms labor markets.

%--- Paper 1 ---
\subsubsection{Skill Content of Technological Change (2003)}

\paragraph{Citation.}
Autor, D.H., Levy, F., \& Murnane, R.J. (2003). ``The Skill Content of Recent
Technological Change: An Empirical Exploration.''
\emph{Quarterly Journal of Economics}, 118(4): 1279-1333.

\paragraph{Core Contribution.}
Computers substitute for routine tasks (both cognitive and manual) but
complement non-routine analytical tasks. This explains why computerization
increases demand for college-educated workers.

\paragraph{Framework Integration.}
\begin{itemize}
  \item \textbf{Task taxonomy:} Routine vs. non-routine, cognitive vs. manual
  \item \textbf{$C^{computer,routine} < 0$:} Computers replace routine tasks
  \item \textbf{$C^{computer,analytical} > 0$:} Computers complement analysis
  \item \textbf{Canonical finding:} Task content predicts employment change
\end{itemize}

\paragraph{Framework Significance.}
This paper (4,000+ citations) established the conceptual foundation for
understanding that technology doesn't simply substitute for ``unskilled''
labor---it substitutes for \emph{routine} tasks at all skill levels.

%--- Paper 2 ---
\subsubsection{Skills, Tasks and Technologies (2011)}

\paragraph{Citation.}
Acemoglu, D. \& Autor, D. (2011). ``Skills, Tasks and Technologies:
Implications for Employment and Earnings.''
\emph{Handbook of Labor Economics}, Vol. 4B, Chapter 12.

\paragraph{Core Contribution.}
Comprehensive theoretical and empirical framework. Tasks as intermediate
inputs produced by workers and machines. Assignment of workers to tasks
is endogenous to technology and wages.

\paragraph{Framework Integration.}
\begin{itemize}
  \item \textbf{Canonical model limitations:} Skill-biased tech change incomplete
  \item \textbf{Task-based extension:} Richer predictions
  \item \textbf{Comparative advantage:} Workers vs. machines by task
  \item \textbf{Wage polarization:} Explains U-shaped changes
\end{itemize}

%--- Paper 3 ---
\subsubsection{Computing Inequality (1998)}

\paragraph{Citation.}
Autor, D.H., Katz, L.F., \& Krueger, A.B. (1998). ``Computing Inequality:
Have Computers Changed the Labor Market?''
\emph{Quarterly Journal of Economics}, 113(4): 1169-1213.

\paragraph{Core Finding.}
Computer use at work strongly associated with higher wages. Industries
with faster computer adoption had faster shifts toward college labor.

\paragraph{Framework Integration.}
\begin{itemize}
  \item \textbf{Empirical foundation:} Computers $\to$ skill demand
  \item \textbf{Industry evidence:} Cross-industry variation
  \item \textbf{Within-industry shift:} Composition change
\end{itemize}

%--- Paper 4 ---
\subsubsection{Trends in U.S. Wage Inequality (2008)}

\paragraph{Citation.}
Autor, D.H., Katz, L.F., \& Kearney, M.S. (2008). ``Trends in U.S.
Wage Inequality: Revising the Revisionists.''
\emph{Review of Economics and Statistics}, 90(2): 300-323.

\paragraph{Core Finding.}
Wage inequality rose throughout 1980s-2000s, not just early 1980s.
Upper-tail inequality (90/50) rose continuously; lower-tail (50/10)
was episodic. Different mechanisms for different parts of distribution.

\paragraph{Framework Integration.}
\begin{itemize}
  \item \textbf{Upper vs. lower tail:} Different dynamics
  \item \textbf{SBTC explains upper tail:} Continuous increase
  \item \textbf{Institutions matter for lower tail:} Minimum wage, unions
\end{itemize}

%--- Paper 5 ---
\subsubsection{Growth of Low-Skill Service Jobs (2013)}

\paragraph{Citation.}
Autor, D.H. \& Dorn, D. (2013). ``The Growth of Low-Skill Service Jobs
and the Polarization of the US Labor Market.''
\emph{American Economic Review}, 103(5): 1553-1597.

\paragraph{Core Finding.}
Employment polarization: growth at top and bottom, decline in middle.
Computerization of routine tasks explains middle-skill decline; rising
demand for non-routine manual service tasks explains low-skill growth.

\paragraph{Framework Integration.}
\begin{itemize}
  \item \textbf{Job polarization:} U-shaped employment growth
  \item \textbf{Routine task decline:} Computerization
  \item \textbf{Service growth:} Non-routine manual tasks resistant
  \item \textbf{Geographic variation:} Commuting zones differ
\end{itemize}

%--- Paper 6 ---
\subsubsection{Why Are There Still So Many Jobs? (2015)}

\paragraph{Citation.}
Autor, D.H. (2015). ``Why Are There Still So Many Jobs? The History and
Future of Workplace Automation.''
\emph{Journal of Economic Perspectives}, 29(3): 3-30.

\paragraph{Core Contribution.}
Accessible synthesis of task framework. Automation substitutes for some
tasks but complements others. Historical perspective shows technology
creates new jobs even as it destroys old ones.

\paragraph{Framework Integration.}
\begin{itemize}
  \item \textbf{Substitution + complementarity:} Both effects
  \item \textbf{Historical evidence:} 200 years of automation
  \item \textbf{Comparative advantage:} Humans vs. machines shifts
  \item \textbf{New work emerges:} Tasks we cannot yet envision
\end{itemize}

\subsection{Part II: The China Shock (5 Papers)}

These papers documented the dramatic local labor market effects of rising
Chinese import competition---the ``China Shock.''

%--- Paper 7 ---
\subsubsection{The China Syndrome (2013)}

\paragraph{Citation.}
Autor, D.H., Dorn, D., \& Hanson, G.H. (2013). ``The China Syndrome:
Local Labor Market Effects of Import Competition in the United States.''
\emph{American Economic Review}, 103(6): 2121-2168.

\paragraph{Core Finding.}
Rising Chinese imports caused large, persistent employment losses in
exposed U.S. commuting zones. Wage decline, labor force exit, increased
transfer payments. Adjustment is slow and incomplete.

\paragraph{Framework Integration.}
\begin{itemize}
  \item \textbf{$\Psi_{trade}$ shock:} China's WTO entry
  \item \textbf{Local effects:} Geographic concentration matters
  \item \textbf{Slow adjustment:} Not instant reallocation
  \item \textbf{$C^{trade,employment} < 0$ locally:} Import competition hurts
\end{itemize}

\paragraph{Framework Significance.}
This paper (3,500+ citations) showed that trade shocks have large,
persistent, geographically concentrated effects---contrary to standard
models predicting rapid adjustment.

%--- Paper 8 ---
\subsubsection{China Shock: Learning from Adjustment (2016)}

\paragraph{Citation.}
Autor, D.H., Dorn, D., \& Hanson, G.H. (2016). ``The China Shock:
Learning from Labor-Market Adjustment to Large Changes in Trade.''
\emph{Annual Review of Economics}, 8: 205-240.

\paragraph{Core Contribution.}
Comprehensive review synthesizing China Shock literature. Adjustment
costs larger and more persistent than previously recognized.

\paragraph{Framework Integration.}
\begin{itemize}
  \item \textbf{Received wisdom overturned:} Adjustment not smooth
  \item \textbf{Full decade+ effects:} Persistence
  \item \textbf{Transfer dependence:} Social safety net absorbs shock
  \item \textbf{Policy implications:} Compensation mechanisms needed
\end{itemize}

%--- Paper 9 ---
\subsubsection{Worker-Level Adjustment (2014)}

\paragraph{Citation.}
Autor, D.H., Dorn, D., Hanson, G.H., \& Song, J. (2014). ``Trade
Adjustment: Worker-Level Evidence.''
\emph{Quarterly Journal of Economics}, 129(4): 1799-1860.

\paragraph{Core Finding.}
Worker-level analysis using Social Security records. Trade shocks reduce
lifetime earnings, increase job mobility, but mobility doesn't restore
earnings. Effects concentrated among older, less-educated workers.

\paragraph{Framework Integration.}
\begin{itemize}
  \item \textbf{Individual trajectories:} Worker-level panel
  \item \textbf{Earnings scarring:} Persistent losses
  \item \textbf{Mobility doesn't help:} New jobs pay less
  \item \textbf{Heterogeneity:} Age, education matter
\end{itemize}

%--- Paper 10 ---
\subsubsection{Import Competition and Marriage Market (2019)}

\paragraph{Citation.}
Autor, D.H., Dorn, D., \& Hanson, G.H. (2019). ``When Work Disappears:
Manufacturing Decline and the Falling Marriage-Market Value of Men.''
\emph{American Economic Review: Insights}, 1(2): 161-178.

\paragraph{Core Finding.}
Trade shocks reduce male employment and earnings, reducing their
``marriageability.'' Result: fewer marriages, more children in
single-parent households, worse child outcomes.

\paragraph{Framework Integration.}
\begin{itemize}
  \item \textbf{Cross-domain spillover:} $C^{trade,family} < 0$
  \item \textbf{Marriage market equilibrium:} Men less ``valuable''
  \item \textbf{Intergenerational effects:} Children affected
  \item \textbf{Framework demonstration:} Economic shocks $\to$ social outcomes
\end{itemize}

%--- Paper 11 ---
\subsubsection{Trade Shocks and Political Polarization (2020)}

\paragraph{Citation.}
Autor, D., Dorn, D., Hanson, G., \& Majlesi, K. (2020). ``Importing
Political Polarization? The Electoral Consequences of Rising Trade Exposure.''
\emph{American Economic Review}, 110(10): 3139-3183.

\paragraph{Core Finding.}
Trade-shocked communities shifted toward ideological extremes: Republicans
in white areas, Democrats in minority areas. Trade contributed to
political polarization and Trump's 2016 victory.

\paragraph{Framework Integration.}
\begin{itemize}
  \item \textbf{$C^{trade,politics}$:} Economic $\to$ political effects
  \item \textbf{Polarization mechanism:} Moderate incumbents lose
  \item \textbf{2016 election:} Trade exposure predicted Trump votes
  \item \textbf{Policy feedback:} Economic policy has political consequences
\end{itemize}

\subsection{Part III: Labor Market Institutions and Policy (5 Papers)}

%--- Paper 12 ---
\subsubsection{Disability Insurance and Employment (2003)}

\paragraph{Citation.}
Autor, D.H. \& Duggan, M.G. (2003). ``The Rise in the Disability Rolls
and the Decline in Unemployment.''
\emph{Quarterly Journal of Economics}, 118(1): 157-205.

\paragraph{Core Finding.}
Disability Insurance (DI) increasingly absorbs displaced workers.
Liberalized screening and higher benefits increased DI rolls,
reducing measured unemployment but also labor supply.

\paragraph{Framework Integration.}
\begin{itemize}
  \item \textbf{Hidden unemployment:} DI as labor force exit
  \item \textbf{Institutional change:} $\Psi_{policy}$ shifted
  \item \textbf{Incentive effects:} Benefits affect behavior
  \item \textbf{Measurement issues:} Official stats misleading
\end{itemize}

%--- Paper 13 ---
\subsubsection{Employment Protection and Productivity (2007)}

\paragraph{Citation.}
Autor, D.H., Kerr, W.R., \& Kugler, A.D. (2007). ``Does Employment
Protection Reduce Productivity? Evidence from U.S. States.''
\emph{Economic Journal}, 117(521): F189-F217.

\paragraph{Core Finding.}
States adopting employment protection (wrongful discharge exceptions)
saw reduced employment flows and productivity growth in affected sectors.

\paragraph{Framework Integration.}
\begin{itemize}
  \item \textbf{$\Psi_{legal}$ affects labor:} Employment protection
  \item \textbf{Trade-off:} Job security vs. dynamism
  \item \textbf{Causal identification:} Cross-state variation
  \item \textbf{Productivity link:} Institutions $\to$ efficiency
\end{itemize}

%--- Paper 14 ---
\subsubsection{Temporary Help Jobs (2010)}

\paragraph{Citation.}
Autor, D.H. \& Houseman, S.N. (2010). ``Do Temporary-Help Jobs Improve
Labor Market Outcomes for Low-Skilled Workers?''
\emph{American Economic Journal: Applied Economics}, 2(3): 96-128.

\paragraph{Core Finding.}
Temporary help placements don't improve long-run employment or earnings
for welfare recipients, despite appearing beneficial in short run.

\paragraph{Framework Integration.}
\begin{itemize}
  \item \textbf{Program evaluation:} Causal effects
  \item \textbf{Selection problem:} Who takes temp jobs?
  \item \textbf{Policy relevance:} Welfare-to-work programs
  \item \textbf{Long-run vs. short-run:} Different conclusions
\end{itemize}

%--- Paper 15 ---
\subsubsection{Outsourcing at Will (2009)}

\paragraph{Citation.}
Autor, D.H. (2009). ``Outsourcing at Will: The Contribution of Unjust
Dismissal Doctrine to the Growth of Employment Outsourcing.''
\emph{Journal of Labor Economics}, 21(1): 1-42.

\paragraph{Core Finding.}
Employment protection laws encouraged firms to outsource workers through
temporary help agencies, shifting legal liability.

\paragraph{Framework Integration.}
\begin{itemize}
  \item \textbf{Institutional arbitrage:} Firms respond to $\Psi_{legal}$
  \item \textbf{Unintended consequences:} Protection $\to$ outsourcing
  \item \textbf{Temp industry growth:} Partly regulatory response
\end{itemize}

%--- Paper 16 ---
\subsubsection{Housing Speculation and Foreclosures (2014)}

\paragraph{Citation.}
Autor, D.H., Palmer, C.J., \& Pathak, P.A. (2014). ``Housing Market
Spillovers: Evidence from the End of Rent Control in Cambridge.''
\emph{Journal of Political Economy}, 122(3): 661-717.

\paragraph{Core Finding.}
Rent decontrol in Cambridge increased property values of never-controlled
units nearby. Positive spillovers from neighborhood quality improvement.

\paragraph{Framework Integration.}
\begin{itemize}
  \item \textbf{$C^{housing,neighborhood} > 0$:} Positive spillovers
  \item \textbf{Policy externalities:} Rent control affected others
  \item \textbf{Natural experiment:} Massachusetts decontrol
\end{itemize}

\subsection{Part IV: AI and the Future of Work (4 Papers)}

%--- Paper 17 ---
\subsubsection{Labor Market Impacts of Technological Change (2022)}

\paragraph{Citation.}
Autor, D.H. (2022). ``The Labor Market Impacts of Technological Change:
From Unbridled Enthusiasm to Qualified Optimism to Vast Uncertainty.''
\emph{NBER Working Paper} No. 30074.

\paragraph{Core Contribution.}
Four paradigms: education race, task polarization, automation-reinstatement,
AI uncertainty. Economic optimism about technology has eroded as
understanding has deepened.

\paragraph{Framework Integration.}
\begin{itemize}
  \item \textbf{Paradigm evolution:} Framework has developed
  \item \textbf{AI different?:} May affect non-routine tasks
  \item \textbf{Uncertainty:} Predictions harder than before
  \item \textbf{Policy urgency:} Need proactive response
\end{itemize}

%--- Paper 18 ---
\subsubsection{New Frontiers: Origins of New Work (2024)}

\paragraph{Citation.}
Autor, D., Chin, C., Salomons, A., \& Seegmiller, B. (2024).
``New Frontiers: The Origins and Content of New Work, 1940-2018.''
\emph{Quarterly Journal of Economics}, forthcoming.

\paragraph{Core Finding.}
60\% of 2018 employment is in job titles that didn't exist in 1940.
New work emerges continuously, creating demand for human labor
even as automation displaces existing tasks.

\paragraph{Framework Integration.}
\begin{itemize}
  \item \textbf{New work creation:} Ongoing process
  \item \textbf{Historical evidence:} 80 years of data
  \item \textbf{Optimism grounded:} Human labor still demanded
  \item \textbf{But:} AI may disrupt pattern
\end{itemize}

%--- Paper 19 ---
\subsubsection{Applying AI to Extend Worker Expertise (2024)}

\paragraph{Citation.}
Autor, D.H. (2024). ``Applying AI to Rebuild Middle Class Jobs.''
\emph{NBER Working Paper} and policy papers.

\paragraph{Core Contribution.}
AI could democratize expertise, enabling workers without extensive
training to perform tasks currently requiring professionals.
Potential to rebuild middle-class jobs, not just destroy them.

\paragraph{Framework Integration.}
\begin{itemize}
  \item \textbf{AI as complement:} $C^{AI,worker} > 0$ possible
  \item \textbf{Expertise democratization:} AI as ``co-pilot''
  \item \textbf{Policy choice:} Direction of technology not fixed
  \item \textbf{Optimistic scenario:} Middle class restored
\end{itemize}

%--- Paper 20 ---
\subsubsection{Concentrating on the Fall of the Labor Share (2020)}

\paragraph{Citation.}
Autor, D., Dorn, D., Katz, L.F., Patterson, C., \& Van Reenen, J. (2020).
``The Fall of the Labor Share and the Rise of Superstar Firms.''
\emph{Quarterly Journal of Economics}, 135(2): 645-709.

\paragraph{Core Finding.}
Labor share declined because ``superstar firms'' (high productivity,
low labor share) gained market share. Reallocation, not within-firm
decline, explains aggregate pattern.

\paragraph{Framework Integration.}
\begin{itemize}
  \item \textbf{Superstar firms:} Winner-take-most dynamics
  \item \textbf{Reallocation effect:} Composition matters
  \item \textbf{Concentration:} Market structure changed
  \item \textbf{Labor share decline:} Not simple automation story
\end{itemize}

\subsection{Synthesis: Autor's Contribution to the Framework}

\paragraph{The Central Insight.}
Autor's work establishes that technology affects labor markets through
\emph{task-specific} complementarities and substitutions:
\begin{equation}
  C^{tech,task_i} \gtrless 0 \quad \text{varies by task type}
\end{equation}

\paragraph{Five Major Contributions.}

\begin{enumerate}
  \item \textbf{Task framework:} Routine/abstract/manual taxonomy
  \item \textbf{Job polarization:} Middle-skill decline explained
  \item \textbf{China Shock:} Trade has large, persistent local effects
  \item \textbf{Slow adjustment:} Labor markets don't instantly clear
  \item \textbf{Cross-domain spillovers:} Economic $\to$ political, family
\end{enumerate}

\paragraph{Framework Elements Established.}

\begin{center}
\renewcommand{\arraystretch}{1.2}
\small
\begin{tabular}{lll}
\hline
\textbf{Element} & \textbf{Papers} & \textbf{Contribution} \\
\hline
Task taxonomy & 1, 2, 6 & Routine/abstract/manual \\
$C^{tech,task}$ varies & 1, 2, 3 & Technology-task complementarity \\
Job polarization & 4, 5 & Middle-skill hollowing out \\
China Shock & 7, 8, 9 & Trade has local effects \\
Cross-domain $C$ & 10, 11 & Economic $\to$ family, politics \\
AI uncertainty & 17, 18, 19 & Future is open \\
\hline
\end{tabular}
\end{center}

\paragraph{Connection to Other Research Programs.}

\begin{center}
\renewcommand{\arraystretch}{1.2}
\begin{tabular}{llllll}
\hline
& \textbf{Fehr} & \textbf{Acemoglu} & \textbf{Shleifer} & \textbf{Heckman} & \textbf{Autor} \\
\hline
Focus & Preferences & Polit. inst. & Legal inst. & Human capital & Labor market \\
$C$ type & Social & Network & Cognitive & Dynamic & Task-based \\
$\Psi$ type & Norms & Institutions & Legal origins & Family & Technology \\
Scale & Individual & Centuries & Decades & Life cycle & Decades \\
\hline
\end{tabular}
\end{center}

\paragraph{Autor + Acemoglu: The Technology Pair.}
Autor and Acemoglu (frequent co-authors) together establish the modern
understanding of technology and labor:
\begin{itemize}
  \item \textbf{Acemoglu:} Directed technical change, automation theory
  \item \textbf{Autor:} Task-based empirics, China Shock
  \item \textbf{Together:} Comprehensive $\Psi_{tech}$ framework
\end{itemize}

\paragraph{The Five Pillars United.}
The $(C, \Psi)$ framework now integrates:
\begin{itemize}
  \item \textbf{Fehr:} Social preferences in context ($\Psi_{norm}$)
  \item \textbf{Acemoglu:} Institutions as cause ($\Psi_{inst}$)
  \item \textbf{Shleifer:} Legal origins and finance ($\Psi_{legal}$)
  \item \textbf{Heckman:} Human capital formation ($C^{skill,investment}$)
  \item \textbf{Autor:} Task-based labor markets ($C^{tech,task}$)
\end{itemize}

Together, these five research programs demonstrate that economic outcomes
depend on context-specific complementarities that the $(C, \Psi)$ framework
is designed to capture.

%%%%%%%%%%%%%%%%%%%%%%%%%%%%%%%%%%%%%%%%%%%%%%%%%%%%%%%%%%%%%%%%%%%%%%%%%%%%%%%
